\subsection{\bt{Amortizing validation and maintenance}{分摊验证与维护开销}}
\label{sec:cost}
\label{sec:opt-eval}
\bi{\emph{Maintenance once, not per agent (Q3).} After the population study's restatement, \system\ repairs each affected definition once, using \PpMavraMaintS\,s of database time per run; private memory repairs per-agent copies in \PpPrivateMaintS\,s. End-to-end maintenance costs \AmMaintBenignS\,s per non-breaking update and \AmMaintRepairS\,s per repaired change. Discarding knowledge on every write costs \ScenRelearnTok\,M relearning tokens per run.}
{\emph{维护只做一次，而不是每个智能体一次（Q3）。}群体实验的数据更正后，\system\ 对每个受影响定义只修复一次，每次运行用 \PpMavraMaintS\ 秒数据库时间；私有记忆逐智能体修复副本，需 \PpPrivateMaintS\ 秒。端到端维护每次正常更新需 \AmMaintBenignS\ 秒、每次修复变化需 \AmMaintRepairS\ 秒。每次写入都丢弃知识，则每次运行需 \ScenRelearnTok\,M token 重新学习。}

\bi{\emph{Shared check results.} Static analysis of all \TpTemplates\ TPC-DS templates yields \TpDefs\ definitions from \TpQueries\ templates, grouping each SELECT block's aggregates by fact table. Updating \code{store\_sales} touches \TpSalesDefs\ definitions but only \TpSalesConds\ distinct conditions. A controlled microbenchmark publishes 4--19 definitions across four families; eight scripted agents use all definitions after each of four updates. Full rechecks cost \CbStagDefOne--\CbStagDefSix\,s as definitions increase, while \system\ stays at \CbCondStagMin--\CbCondStagMax\,s, or \CbStagCondSix\% of full-recheck time at 19 definitions. A general cache keyed by validation SQL and table versions costs \CbStagCacheSix\% in the staggered run; with simultaneous arrivals, it costs \CbBurstCacheSix\% versus \CbBurstCondSix\% for \system. Condition identity also recognizes equivalent key orders that produce different SQL. Without shared conditions, all methods cost \CbZeroMin--\CbZeroMax\% of full rechecking. These controlled microbenchmarks isolate the contribution of condition overlap to maintenance savings; the ratios remain low at 4 million rows (\CbFourStagCond\%), with 32 agents (\CbAgentsStagCond\%), and up to 16 million rows (\EoneRatioSixteen--\EoneRatioFour\%).}
{\emph{共享的检查结果。}对全部 \TpTemplates\ 个 TPC-DS 模板作静态分析，将每个 SELECT 块的聚合按事实表归组，从 \TpQueries\ 个模板中导出 \TpDefs\ 个定义。更新 \code{store\_sales} 触及 \TpSalesDefs\ 个定义，但只有 \TpSalesConds\ 个不同条件。受控微基准在四个族中发布 4--19 个定义，四次更新后各让八个脚本智能体使用全部定义。定义数增加时，整定义重查耗时从 \CbStagDefOne\ 增至 \CbStagDefSix\ 秒，\system\ 保持在 \CbCondStagMin--\CbCondStagMax\ 秒，19 个定义时为整定义重查的 \CbStagCondSix\%。以验证 SQL 与表版本为键的通用缓存，在错峰运行中为 \CbStagCacheSix\%；同时到达时为 \CbBurstCacheSix\%，\system\ 为 \CbBurstCondSix\%。条件标识还能识别生成不同 SQL 的等价键顺序。没有共享条件时，各方法都需整定义重查的 \CbZeroMin--\CbZeroMax\%。这些受控微基准分离出条件重叠对维护节省的贡献；重查耗时比例在 400 万行（\CbFourStagCond\%）、32 个智能体（\CbAgentsStagCond\%）以及直到 1600 万行（\EoneRatioSixteen--\EoneRatioFour\%）时保持较低。}

\bi{\emph{One optimization speeds up every agent.} We evaluate an optimization pass over the \OptDefs\ learned definitions on the one-million-row retail workload. \OptPublished\ of the \OptDefs\ published definitions received a key-range revision that reduced query execution time by \OptSaveMin--\OptSaveMax\% relative to the previous revision (mean \OptSaveMean\%, every 95\% interval below zero), with equal results on \OptPeriods\ periods each, and \OptHoldAdopted\ of the \OptHoldTasks\ held-out queries ran the faster form because agents copy the SQL they retrieve. A separate system-level test publishes and optimizes \PbDefs\ definitions (store revenue, return amount, return rate, and electronics revenue), then re-keys one month of the date dimension so that only date-key contiguity fails: maintenance invalidated all \PbDefs\ key-range revisions in this test and restored their previous revisions within \PbFallbackMsMax\,ms each, and all \PbAnswers\ answers served afterwards are correct, whereas the key-range SQL would have been wrong by up to \PbStaleErrMax\%.}
{\emph{一次优化让每个智能体都更快。}我们在 100 万行销售的零售负载上，针对学得的 \OptDefs\ 个定义评估一轮优化。已发布的 \OptDefs\ 个定义中有 \OptPublished\ 个得到键区间版本，相对前一版本使查询执行耗时减少 \OptSaveMin--\OptSaveMax\%（平均减少 \OptSaveMean\%，每个 95\% 区间都低于零），每个都在 \OptPeriods\ 个期间上结果相同；由于智能体照抄检索到的 SQL，\OptHoldTasks\ 条留出查询中有 \OptHoldAdopted\ 条用了更快的写法。另一项独立的系统层测试发布并优化 \PbDefs\ 个定义（门店营业额、退货金额、退货率与电子品类营业额），随后给日期维度的一个月重新编键，使只有日期键连续性失败：维护使该测试中全部 \PbDefs\ 个键区间版本失效，并在每个 \PbFallbackMsMax\ 毫秒内恢复其前一版本，之后给出的 \PbAnswers\ 个答案全部正确，而键区间 SQL 的误差会高达 \PbStaleErrMax\%。}
